\begin{tikzpicture}[
  font=\scriptsize,
  sistem/.style={draw, rounded corners, fill=blue!8, text width=1.6cm,
                 minimum height=0.55cm, align=center, font=\scriptsize},
  tujuan/.style={draw, rounded corners, fill=orange!12, text width=1.75cm,
                 minimum height=0.55cm, align=center, font=\scriptsize},
  panah/.style={-{Stealth[length=2mm]}, thick}
]
  % Kiri: setiap SIMRS disambung langsung ke setiap tujuan.
  \node[font=\scriptsize\bfseries] at (1.7,2.65) {Point-to-point saat ini};
  \foreach \nama [count=\nomor] in {Jakarta, Bekasi, Depok, Bogor, Bandung} {
    \node[sistem] (a\nomor) at (0,{2.6-\nomor*0.85}) {SIMRS \nama};
  }
  \foreach \nama [count=\nomor] in {BPJS Kesehatan, SATUSEHAT, Medikara+,
                                    ERP keuangan} {
    \node[tujuan] (t\nomor) at (3.4,{1.75-\nomor*0.85}) {\nama};
  }
  \foreach \asal in {1,...,5} {
    \foreach \tuju in {1,...,4} {
      \draw[red!55] (a\asal.east) -- (t\tuju.west);
    }
  }
  \node[font=\tiny, align=center, anchor=north] at (1.7,-2.05)
    {5 $\times$ 4 = 20 sambungan,\\tiap sambungan dibangun dan dirawat sendiri};

  \draw[gray, dashed] (4.6,2.85) -- (4.6,-2.75);

  % Kanan: setiap sistem cukup tersambung ke satu lapisan bersama.
  \node[font=\scriptsize\bfseries] at (7.55,2.65) {Lewat lapisan API bersama};
  \foreach \nama [count=\nomor] in {Jakarta, Bekasi, Depok, Bogor, Bandung} {
    \node[sistem] (b\nomor) at (5.7,{2.6-\nomor*0.85}) {SIMRS \nama};
  }
  \node[draw, rounded corners, fill=green!10, text width=0.95cm,
        minimum height=4.1cm, align=center, font=\scriptsize\bfseries]
    (lapisan) at (7.55,0.05) {Lapisan API bersama};
  \foreach \nama [count=\nomor] in {BPJS Kesehatan, SATUSEHAT, Medikara+,
                                    ERP keuangan} {
    \node[tujuan] (u\nomor) at (9.45,{1.75-\nomor*0.85}) {\nama};
  }
  \foreach \asal in {1,...,5} {
    \draw[praditagreen, thick] (b\asal.east) -- (b\asal.east -| lapisan.west);
  }
  \foreach \tuju in {1,...,4} {
    \draw[praditagreen, thick] (u\tuju.west -| lapisan.east) -- (u\tuju.west);
  }
  \node[font=\tiny, align=center, anchor=north] at (7.55,-2.05)
    {5 + 4 = 9 sambungan\\ke satu tempat yang dapat dikelola};
\end{tikzpicture}
